% 译者说明（Lead）
\chapter*{译者说明}
\addcontentsline{toc}{chapter}{译者说明}
\markboth{译者说明}{译者说明}

本译本依据岩波书店日译本《K理論》（M. F. Atiyah 著，松尾信一郎 监译，川边治之 译）译出；该日译本又据 Atiyah 原著 \textit{K-Theory}（W. A. Benjamin, 1967 年初版）译成。原著由 1964 年秋在哈佛大学的讲义整理而成，另收两篇已发表论文作为附录 B、C。

\medskip
\noindent\textbf{体例说明}
\begin{itemize}[leftmargin=2em,itemsep=2pt]
\item 原书页码以注释形式标在每页起始处（源码中为 \verb|% [原文 p.N]|）；PDF 页号 = 原书页码 + 9。卷首「引言」用罗马页码（p.v）。
\item 公式编号、定义／定理／命题／引理／推论等陈述编号一律照录原书，不重排、不改号。原书公式按「章.节.序号」编号（如式 (2.4.2)），陈述编号同此（如定义 1.4.3、命题 B.2.5）。
\item 原书正文的笔误照原样排，只在边注中（源码中为 \verb|% [Tutor Note: ...]|）指出并给出正确形式；批注统一收在书末的《译校批注》里。
\item 术语以书末《术语对照表》为准；人名、模型名、缩写保留英文原形。日译本的译者注（页脚「＊（訳注）」）照样译出，注为脚注。
\item 每节前有「本节重点」，每章（含附录）末有「本章回顾与评注」，均为译者所加，用以提示结构与易错点，可在导言区用 \verb|\showtakeawayfalse|、\verb|\showchapterreviewfalse| 关闭。
\item 附录 B、C 是 Atiyah 两篇独立论文的重印，文风与讲义部分略有差别；两处末尾的「参考文献」照原书抄录（条目为西文原文，不翻译）。
\item 原书目录（p.vii–viii）不重排，由本书自动生成的目录代替；原书索引（p.203–207）未收，见 README 的「已知限制」。
\end{itemize}

\medskip
\noindent 译文力求准确，但错误在所难免，恳请读者指正。
